% !TeX program = xelatex
% 本文件由 assemble.py 自 OI-Docs 原稿抽取生成，请勿手改（会被覆盖）。
% 源文件：Number Theory/number_theory_teacher.tex
% 源行号：3762-4149 / 4152-4350 / 4353-4639（数论）+ 1642-1728（博弈）
% 说明：已删去独立成册所需的导言区与页码重置；章节正文与交叉引用原样保留。

% ==========================================================================
% 附录：A--D。章号由 main.tex 中的 \appendix 切换为 A、B、C、D。
% ==========================================================================

\chapter{分层训练与参考解析}
\begin{chapterintro}[章节导读]
前十五章建立的是一条从基本数论工具一路走到高级专题的主线；从这里开始，正文不再继续引入新的知识体系，而是转向训练与教学资源。附录 A 把前面学过的方法重新组合进分层练习中，重点从“会不会公式”转向“能不能选对工具、检查条件并写出完整推导”。
\end{chapterintro}

本章是新增训练，不替代前文原题。学生版只保留题目；教师版在每题后给出结果、关键步骤与评价要点。建议先独立完成，再检查条件和推导，而不是仅比对最后一个数。

\section{A 组：整除、通解与基础建模}
\begin{problem}{A1：从 gcd 回代到不定方程}
求 $\gcd(414,662)$，给出 $414x+662y=6$ 的一组整数解与通解。
\end{problem}

\begin{pitfall}[讲解：先看这题到底在做什么]
先把 gcd 的回代结果写成一个线性组合，再整体乘到右端。这里练的是“exgcd 的输出怎么变成方程特解”，而不是重新证明裴蜀定理。
\end{pitfall}

\sol gcd 为 $2$。欧几里得回代得 $2=8\cdot414-5\cdot662$，乘三后得到 $(x_0,y_0)=(24,-15)$。通解为
\[
 x=24+331t,\qquad y=-15-207t,\quad t\in\Z.
\]

\begin{tip}[三档练习的用法]
A 组只用到第 1--3 章，做完应该能独立写出通解区间；B 组对应第 4--6 章，重点是\textbf{先列周期再套公式}；C、D 组要求你识别工具而不是回忆公式。建议每题限时 $15$ 分钟，写不出来就读解析的第一句，仍然卡住就先跳到下一组，回头再战。
\end{tip}

评价时应分别检查：整除条件、缩放方向、通解步长；只写一个特解不能算完成。


\begin{problem}{A2：正整数解的计数}
求 $14x+21y=147$ 的全部正整数解、解数以及两个变量的极值。
\end{problem}

\begin{pitfall}[讲解：先看这题到底在做什么]
先写出全部整数解，再把 $x>0,y>0$ 各自产生的条件变成同一个参数 $t$ 的上下界。这样“数多少组正整数解”就完全变成了整数区间计数。
\end{pitfall}

\sol 化为 $2x+3y=21$。取 $x=3t,y=7-2t$，约束给 $1\le t\le3$，所以解为 $(3,5),(6,3),(9,1)$。解数 $3$，$x$ 最小 $3$、最大 $9$，$y$ 最小 $1$、最大 $5$。强调两个最小值不来自同一组解。


\begin{problem}{A3：约分会改变模数}
解 $18x\equiv30\pmod{42}$，并在区间 $[0,42)$ 中列出全部解。
\end{problem}

\begin{pitfall}[讲解：先看这题到底在做什么]
同余式两边约掉公共因子以后，模数也必须同步除掉这个因子。先得到约分后的最小模数，再列出全部同余解；只写一个代表元会漏掉区间里的其他解。
\end{pitfall}

\sol 除 gcd $6$ 得 $3x\equiv5\pmod7$，故 $x\equiv4\pmod7$。区间内解为 $4,11,18,25,32,39$，共有六个。只给 $4$ 是漏解；把约分后模数仍写成 $42$ 是错误。


\begin{problem}{A4：CRT 与下界}
求满足 $x\equiv1\pmod4$、$x\equiv3\pmod6$ 且 $x\ge40$ 的最小整数。
\end{problem}

\begin{pitfall}[讲解：先看这题到底在做什么]
CRT 负责找到所有满足同余条件的数，即一个同余类；题目还有“至少多大”的额外限制，所以最后还要从这个同余类中找第一个不小于下界的代表元。
\end{pitfall}

\sol 令 $x=1+4t$，$4t\equiv2\pmod6$，故 $t\equiv2\pmod3$，合并为 $x\equiv9\pmod{12}$。不小于 $40$ 的首项为 $45$。直接输出最小非负代表元 $9$ 忽略了大小约束。


\begin{problem}{A5：一个质因数分解对应多种函数}
计算 $\varphi(72),\mu(72),\tau(72)$ 与 $\sum_{d\mid72}\varphi(d)$。
\end{problem}

\begin{pitfall}[讲解：先看这题到底在做什么]
把同一个质因数分解同时喂给 $\varphi,\mu,	au$，可以很好地训练“各函数到底看分解的什么信息”：$\varphi$ 看质因子及其幂次，$\mu$ 只关心是否出现平方因子，$	au$ 直接把各指数加一后相乘。
\end{pitfall}

\sol $72=2^3\cdot3^2$，所以 $\varphi(72)=72(1-1/2)(1-1/3)=24$，$\mu(72)=0$，$\tau(72)=(3+1)(2+1)=12$。最后一项直接由欧拉恒等式等于 $72$，无需枚举十二个约数。评价重点是区分“有平方因子”与“不同素因子个数”。


\begin{problem}{A6：单次逆元与批量逆元}
求 $37^{-1}\pmod{101}$；再说明模 $17$ 同时求 $2,3,5$ 的逆元时怎样只做一次一般求逆。
\end{problem}

\begin{pitfall}[讲解：先看这题到底在做什么]
先用一次 exgcd 得到基准逆元，再观察前缀积之间的乘积逆元关系。批量法的节省之处不在于“exgcd 更快”，而在于原本很多次求逆被一次逆元和若干乘法替代。
\end{pitfall}

\sol $1=11\cdot101-30\cdot37$，故逆元为 $71$。批量问题先算前缀积 $1,2,6,13$，求 $13^{-1}=4$，再逆推得到 $5^{-1}=7,3^{-1}=6,2^{-1}=9$。每一步应用的是乘积逆元关系，而不是将三个独立 exgcd 称为“批量”。


\begin{problem}{A7：环的连通分量}
在 $18$ 个点的环上每次加 $12$，有几个循环？从点 $3$ 出发能否到点 $4$？列出点 $3$ 所在循环。
\end{problem}

\begin{pitfall}[讲解：先看这题到底在做什么]
每走一步就是给编号加 $k$ 并模 $n$。能走到哪些点只由 $\gcd(n,k)$ 决定；因此先算 gcd 就能直接知道环被分成多少个分量，而不是逐点模拟。
\end{pitfall}

\sol gcd 为 $6$，所以有六个循环，每个长三。每次保持模 $6$ 余数，$3$ 与 $4$ 不在同一循环。对应循环为 $3\to15\to9\to3$。可以进一步追问：到某个点的最少步数与“是否同环”不是同一个问题。


\begin{problem}{A8：无整数解与无正整数解}
比较 $6x+10y=7$ 与 $6x+10y=2$。第二个方程中，单独可取得的最小正 $x$、最小正 $y$ 分别是多少？它们能否同时取得？
\end{problem}

\begin{pitfall}[讲解：先看这题到底在做什么]
先区分两个问题：$d
mid c$ 时连整数解都没有；$d\mid c$ 时整数解存在，但正整数限制可能让参数区间为空。即使没有同时为正的解，最小正 $x$ 和最小正 $y$ 仍可能分别存在。
\end{pitfall}

\sol 前者因为 $2\nmid7$ 而无整数解。后者有解 $(-3,2)$，通解 $x=-3+5t,y=2-3t$，没有同时为正的参数。最小正 $x=2$（对应 $y=-1$），最小正 $y=2$（对应 $x=-3$），不能同时取得。这正是 P5656 特殊输出分支的含义。


\section{B 组：指数、赋值与组合数}
\begin{problem}{B1：欧拉指数与真实周期}
求 $7^{123456789}\bmod40$，并分别使用欧拉函数与阶解释降幂依据。
\end{problem}

\begin{pitfall}[讲解：先看这题到底在做什么]
欧拉定理给出的 $\varphi(m)$ 只是一个一定有效的指数，不一定是最小周期。先通过小次幂找到真实的阶，再决定指数应该对什么数取模，能避免把“周期”和“可用上界”混为一谈。
\end{pitfall}

\sol $\gcd(7,40)=1$，$\varphi(40)=16$，指数模 $16$ 为 $5$；而 $7^2\equiv9,7^4\equiv1$，真实阶为 $4$。指数模 $4$ 为一，答案 $7$。欧拉定理给出一个周期，未必给最小周期。


\begin{problem}{B2：非互质大指数}
求 $6^{10^{100}}\bmod72$，并说明不计算任何大指数即可得到答案的原因。
\end{problem}

\begin{pitfall}[讲解：先看这题到底在做什么]
这里不要机械套欧拉降幂。直接看各素因子的赋值什么时候达到模数所需的指数，一旦整个数已经被模数整除，答案就立即是 $0$。这是非互质底数时最直接的判定方式。
\end{pitfall}

\sol $72=2^3\cdot3^2$，$6^e$ 中两个素数的赋值均为 $e$。当 $e\ge3$ 即同时被两个素数幂整除，故答案为零。此处直接按赋值判零，比构造欧拉链更简单。不要把所有巨大指数问题都机械套用递归降幂。


\begin{problem}{B3：先列周期，再写离散对数通解}
求 $\ord_{17}(4)$，并解 $4^x\equiv13\pmod{17}$。
\end{problem}

\begin{pitfall}[讲解：先看这题到底在做什么]
先把 $4$ 的幂循环写出来，阶自然就出现了。找到最小解以后，所有解都只是在这个周期上平移，因此通解就是“最小解 + 阶的倍数”。
\end{pitfall}

\sol $4^2\equiv16=-1,4^4\equiv1$，所以阶为四。周期为 $1,4,16,13$，最小非负解为三，全部非负解为 $x=3+4t$（$t\ge0$）。这里模 $17$ 的阶不是 $\varphi(17)=16$。


\begin{problem}{B4：离散对数的有限前缀}
分别求 $6^x\equiv12\pmod{18}$ 与 $6^x\equiv0\pmod{18}$ 的最小非负解。
\end{problem}

\begin{pitfall}[讲解：先看这题到底在做什么]
底数与模数不互质时，序列可能先快速坍缩，再进入一个短周期。此时不能定义一个正常的乘法阶，因此要先观察前缀行为，再判断是否存在目标值。
\end{pitfall}

\sol 序列为 $1,6,0,0,\ldots$，第一式无解，第二式最小解为二。这一序列不是从指数零开始的纯周期；也不存在 $\ord_{18}(6)$。学生若写“指数对阶取模”，应先让其指出阶为何不存在。


\begin{problem}{B5：奇素数 LTE}
求 $\nu_3(10^{81}-1)$，并列出所使用的全部前提。
\end{problem}

\begin{pitfall}[讲解：先看这题到底在做什么]
LTE 不是只背最后一个公式。讲解时先逐条检查素数、整除关系和指数条件，再把赋值直接代入。学生能明确说出“为什么这条 LTE 能用”，比记住结果更重要。
\end{pitfall}

\sol 素数三为奇数，$3\mid10-1$，$3\nmid10$、$3\nmid1$，指数为正。故结果为 $\nu_3(9)+\nu_3(81)=2+4=6$。满分答案必须写出这些条件，而不是只写六。


\begin{problem}{B6：二的赋值}
求 $\nu_2(5^{12}-1)$，用一般偶指数形式与 $4\mid x-y$ 的简化形式各算一次。
\end{problem}

\begin{pitfall}[讲解：先看这题到底在做什么]
模 $2$ 的 LTE 与奇素数版本不一样，最稳妥的做法是先写通式，再看额外条件能否简化。最后两种形式得到相同结果时，可以把它当作一次条件自检。
\end{pitfall}

\sol 一般式为 $\nu_2(4)+\nu_2(6)+\nu_2(12)-1=2+1+2-1=4$。由于 $4\mid5-1$，简化式为 $\nu_2(4)+\nu_2(12)=4$。两种形式一致是一次很好的条件检查。


\begin{problem}{B7：阶乘赋值与组合数赋值}
求 $\nu_3(100!)$ 与 $\nu_2\binom{100}{50}$。
\end{problem}

\begin{pitfall}[讲解：先看这题到底在做什么]
先把阶乘中的 $p$ 因子数量分别算出来，再做减法得到组合数的赋值。Legendre 公式和二进制数位和公式是两种不同视角，适合让学生验证同一个结果。
\end{pitfall}

\sol 第一项为 $33+11+3+1=48$。第二项为 $\nu_2(100!)-2\nu_2(50!)=97-2\cdot47=3$。也可用二进制数位和：$100$、$50$ 都有三个一，所以 $3+3-3=3$。结果说明该组合数能被八整除，不能被十六整除。


\begin{problem}{B8：一次 Lucas 计算}
求 $\binom{20}{6}\bmod7$。若把下标改成八，结果是否仍非零？
\end{problem}

\begin{pitfall}[讲解：先看这题到底在做什么]
Lucas 定理的关键是把 $n,k$ 写成 $p$ 进制，然后逐位取组合数。只要出现某一位的下标超过上标，对应因子就是 $0$；没有越界时才继续相乘。
\end{pitfall}

\sol $20=(26)_7,6=(06)_7$，余数为 $\binom20\binom66=1$。$8=(11)_7$，余数为 $\binom21\binom61=12\equiv5$，仍非零。这两题都没有位上越界；改成 $7$ 进制某位超过上标的数才会自动为零，不能仅凭下标“变大”判断。


\begin{problem}{B9：合数模数不能直接除阶乘}
求 $\binom{12}{4}\bmod18$，要求按模二和模九分别解释。
\end{problem}

\begin{pitfall}[讲解：先看这题到底在做什么]
组合数取模时最危险的地方就是把普通除法机械搬进模环。先分别处理素数幂中的 $p$ 因子和单位部分，再用 CRT 合并；尤其要明确“分母没有逆元”时不能直接除。
\end{pitfall}

\sol 模二：$\nu_2=10-3-7=0$，组合数为奇数，余数一。模九：$\nu_3=5-1-2=2$，至少含 $3^2$，余数零。CRT 合并“奇数且为九的倍数”，模十八余九。实际 $\binom{12}{4}=495$。不允许把 $4!\bmod18=6$ 的“逆元”代入，因为它根本不存在。


\begin{problem}{B10：辨认计数对象}
(1) P7488 的三角形模型中 $n=4,m=2$ 的组数是多少？(2) 五件不同礼物分别给两位有编号的人两件、一件，有多少方案？
\end{problem}

\begin{pitfall}[讲解：先看这题到底在做什么]
两道题表面都是“选一些对象”，但排列是否算不同、接收者是否有身份，决定了计数公式。先明确对象是什么、哪些选择不能重复，再决定是否需要乘或除阶乘。
\end{pitfall}

\sol (1) $\binom42^2\cdot2!=72$，也可由 $R(3,2)+(8-2)R(3,1)=18+6\cdot9=72$。 (2) $\binom52\binom31=30$。第一题三角形组无序，第二题接收者有身份，但两题最终都通过恰好计数避免重复；不能见到“组”就随意除阶乘。


\section{C 组：反演与亚线性求和}
\begin{problem}{C1：手算一次杜教筛}
计算 $S_\varphi(12)$ 和 $S_\mu(12)$，并写出前者按商分块后的递归式。
\end{problem}

\begin{pitfall}[讲解：先看这题到底在做什么]
先把递推式中的每一个块看成一个整体，而不是逐项展开。$\floor{n/j}$ 相同的连续 $j$ 会共享同一个状态，所以真正要算的是块长和对应的递归状态；手算一次后再看代码会清楚很多。
\end{pitfall}

\sol
\[
S_\varphi(12)=78-S_\varphi(6)-S_\varphi(4)-S_\varphi(3)-2S_\varphi(2)-6S_\varphi(1)=46.
\]
$S_\mu(10)=-1,\mu(11)=-1,\mu(12)=0$，所以 $S_\mu(12)=-2$。本题检查的是块长与从 $j=2$ 开始，漏掉这两处任意一处都会出错。


\begin{problem}{C2：矩形上的 gcd 与互质计数}
求 $\sum_{i\le4,j\le6}\gcd(i,j)$，以及其中互质有序数对的个数。
\end{problem}

\begin{pitfall}[讲解：先看这题到底在做什么]
这两种计数都能拆成“公共因子 $d$”，但权重分别来自 $\varphi$ 和 $\mu$。先把计数对象说清楚，再写权重，能避免把两个看起来相似的公式混在一起。
\end{pitfall}

\sol gcd 和为 $24+6+4+2=36$；互质数对为 $24-6-2=16$。前者权重是 $\varphi(d)$，后者是 $\mu(d)$，两式商因子相同，不能把权重混用。也不能在矩形问题中将两个商都写成 $\floor{4/d}$。


\begin{problem}{C3：gcd 的平方应该使用哪个系数？}
推导 $\sum_{i,j\le n}\gcd(i,j)^2$ 的一维求和式，给出系数在素数幂处的值，并计算 $n=3$ 的答案。
\end{problem}

\begin{pitfall}[讲解：先看这题到底在做什么]
平方 gcd 对应的是 Jordan 型卷积系数，而不是把 $\varphi$ 直接平方。最稳妥的解释是从 $\id_2=J_2*\one$ 出发，再用同样的公共因子计数套路。
\end{pitfall}

\sol 令 $J_2=\id_2*\mu$，则 $\id_2=J_2*\one$，所以答案是
\[
 \sum_{d\le n}J_2(d)\floor{n/d}^2,\qquad
 J_2(p^e)=p^{2e}-p^{2(e-1)}.
\]
$n=3$ 时为 $9\cdot1+1\cdot3+1\cdot8=20$。系数不是 $\varphi(d)^2$，因为卷积不保持这种逐点平方关系。


\begin{problem}{C4：PN 思想的最简单计数应用}
求不超过三十的平方自由数个数，要求不逐个判定。
\end{problem}

\begin{pitfall}[讲解：先看这题到底在做什么]
这个例子只是在数没有平方因子的整数，但它完整体现了 PN 的套路：先做局部卷积消去主要贡献，再留下平方数支撑，最后只枚举 $d\le\sqrt n$。
\end{pitfall}

\sol 使用 $\sum_{d\le\sqrt{30}}\mu(d)\floor{30/d^2}$，得到 $30-7-3+0-1=19$。讲解时可让学生说明为什么含多个平方因子的数会通过莫比乌斯系数正确容斥。


\begin{problem}{C5：PN 卷积商的交叉项}
设积性函数满足 $f(p^k)=p^{2k}+[k\ge2]$。取 $g(n)=n^2$ 并令 $f=g*h$，求 $h(p),h(p^2),h(p^3)$，设计 $S_f(n)$ 的稀疏求和式。
\end{problem}

\begin{pitfall}[讲解：先看这题到底在做什么]
卷积商在 $p^3$ 处会混入低次项的贡献，所以不能直接写成“原函数减基准函数”。先按素数幂递推把前面的项扣干净，再看剩余支撑为什么仍然稀疏。
\end{pitfall}

\sol $h(p)=0$，$h(p^2)=1$，而
\[
 h(p^3)=f(p^3)-g(p^3)-g(p)h(p^2)=1-p^2.
\]
故不能对所有幂次都写成 $f(p^k)-g(p^k)$。有
\[
 S_f(n)=\sum_{d\in\mathrm{PN},d\le n}h(d)\,
 \frac{q(q+1)(2q+1)}6,\quad q=\floor{n/d}.
\]
在局部系数及枚举准备完成后，只需 $O(\sqrt n)$ 个支撑项；若取模且分母不可逆，要整数约分后再取模。


\begin{problem}{C6：验证 P4240 的简化系数}
计算 $f(6),f(12)$，其中 $f(k)=\sum_{d\mid k}d\mu(k/d)/\varphi(d)$；再用 $G$ 形式算 $\sum_{i,j\le3}\varphi(ij)$。
\end{problem}

\begin{pitfall}[讲解：先看这题到底在做什么]
这一题适合专门检查“局部系数到底有没有算对”。先在整数或有理数层面验证几个小 $k$，确认系数结构正确，再把它们放回模意义下的实现。
\end{pitfall}

\sol $6$ 平方自由，$f(6)=1/\varphi(6)=1/2$；$12$ 含平方因子，$f(12)=0$。后一个求和中，$k=1,2,3$ 的贡献为 $16,1,2$，合计 $19$。分数可先在有理数中理解，再在题目模数中用逆元表示。


\section{D 组：证明辨析与综合拓展}
\begin{problem}{D1：非平凡平方根为何能分解？}
列出 $x^2\equiv1\pmod{15}$ 的所有解。用其中一个非平凡解找到十五的两个素因子。
\end{problem}

\begin{pitfall}[讲解：先看这题到底在做什么]
CRT 告诉我们模素因子分别取 $1$ 或 $-1$，组合起来可以得到非平凡平方根。只要它既不是 $1$ 也不是 $-1$，那么与 $V-1,V+1$ 的 gcd 就会分别捞出不同因子。
\end{pitfall}

\sol CRT 分别在模三、模五选择 $\pm1$，得 $1,4,11,14$。选 $V=4$，则 $\gcd(V-1,15)=3$，$\gcd(V+1,15)=5$。若选 $V=14=-1$，后一 gcd 等于十五，不会分裂；这正是随机分解步骤为何排除 $-1$。


\begin{problem}{D2：相同阶不代表同一个子群}
比较模十五的 $2,7$ 的阶。是否存在 $a\ge0$ 使 $2^a\equiv7\pmod{15}$？说明 P5605 中原根条件不可删除。
\end{problem}

\begin{pitfall}[讲解：先看这题到底在做什么]
“阶相同”和“生成同一个子群”是两件事。先列出两个元素的幂集合，再直接比较集合是否相同，就能看到为什么在非循环单位群里不能只靠阶判断子群包含。
\end{pitfall}

\sol 两者阶均为四，但 $2$ 的幂集合为 $\{1,2,4,8\}$，$7$ 的幂集合为 $\{1,7,4,13\}$，后者并不属于前者生成的子群。模十五的单位群不循环；阶的整除关系不能替代子群包含关系。


\begin{problem}{D3：一步边也可能产生多个路径余数}
两个点仅通过一条边权二的无向边相连，模数为七。允许反复走边，从一端到另一端能得到哪些权值余数？
\end{problem}

\begin{pitfall}[讲解：先看这题到底在做什么]
先把路径权值写成一个几何级数，再研究 $2^h$ 在允许的奇数 $h$ 上能出现哪些余数。这里必须同时关注图上的路径结构和数值状态，不能只看其中一个。
\end{pitfall}

\sol 路径长度 $h$ 必须为奇数，原权值为 $2(1+2+\cdots+2^{h-1})=2(2^h-1)$。$2^h$ 在奇数 $h$ 上可取 $2,1,4$，所以余数为 $2,0,6$。图的拓扑连通性与数值状态连通性必须同时考虑。


\begin{problem}{D4：二维卷积的一次完全消去}
取 $S(n)=\tau(n)$。证明
\[
 \tau(xy)=\sum_{d\mid\gcd(x,y)}\mu(d)\tau(x/d)\tau(y/d),
\]
并写出 $\sum_{i,j\le n}\tau(ij)$ 的一维表达式。
\end{problem}

\begin{pitfall}[讲解：先看这题到底在做什么]
这题真正要看的不是最后那一行公式，而是局部生成函数里的 $1-XY$。它意味着所有坐标轴上的主要贡献都被消掉，只剩对角线上的稀疏支撑，于是二维求和又降成一维。
\end{pitfall}

\sol 固定素数，$S(p^k)=k+1$。二元局部级数为
\[
 \sum_{a,b\ge0}(a+b+1)X^aY^b
 =\frac{1-XY}{(1-X)^2(1-Y)^2}.
\]
分母对应 $\tau(x)\tau(y)$，剩余 $1-XY$ 表示卷积商只在局部 $(p,p)$ 为 $-1$，因此全局只在 $x=y=d$ 为 $\mu(d)$。答案为
\[
 \sum_{d\le n}\mu(d)S_\tau(\floor{n/d})^2.
\]
$n=3$ 时为 $25-1-1=23$。这展示了“算法二”在某些函数上把所有非对角残差完全消掉，而不是只得到一个较稀疏的支撑。


\begin{problem}{D5：从最大素因子开始还原}
已知 $M=1080=N\varphi(N)$ 且解存在，求 $N$，并解释为什么不是直接取 $\sqrt M$。
\end{problem}

\begin{pitfall}[讲解：先看这题到底在做什么]
已知 $N\varphi(N)$ 后，可以从最大的素因子开始拆，因为 $N$ 和 $\varphi(N)$ 中各素因子的指数关系很受限制。每还原出一个素因子，就更新剩余量，再继续处理下一层。
\end{pitfall}

\sol $1080=2^3\cdot3^3\cdot5$，最大素因子五的指数为一，对应 $N$ 中一个五；除去 $5\cdot4=20$ 后剩五十四。三的指数为三，对应 $N$ 中两个三；除去 $3^3\cdot2=54$ 后结束。所以 $N=45$，$\varphi(45)=24$，乘积为 $1080$。$\varphi(N)$ 一般不等于 $N$，平方根只能提供量级，不能给答案。


\begin{problem}{D6：构造题还需要全局去重}
当 $k=2$，从种子 $(1,2,6)$ 写出两个孩子。为什么它们都不能作为该种子之后的第二组输出？沿 BFS 继续找一组可用输出。
\end{problem}

\begin{pitfall}[讲解：先看这题到底在做什么]
构造出一个满足方程的三元组只是局部正确；题目通常还要求不同组之间不能重复使用整数。因此搜索状态必须同时记录“这个组是否合法”和“这些数是否已经在全局出现过”。
\end{pitfall}

\sol 两个孩子是 $(2,6,15)$、$(1,6,12)$，都与种子共享已使用整数。不能因为“三元组不相同”就输出。按正文孩子顺序 BFS，继续可找到 $(15,40,104)$；它与 $1,2,6$ 无交，且
\[
15^2+40^2+104^2=12641=2(15\cdot40+40\cdot104+15\cdot104)+1.
\]
把“单组合法”和“所有组之间合法”分别验证，是构造题的通用习惯。





\chapter{核心模板与接口}
\begin{chapterintro}[章节导读]
附录 A 训练的是分析与推导，接下来需要把这些结论真正落到程序接口上。本附录集中整理高频数论模板，并把输入前提、溢出边界、返回值以及空间成本写在代码旁边，方便把数学证明转化为可靠实现。
\end{chapterintro}

\section{模板使用说明}

\begin{tip}[附录代码怎么读]
先读注释里的 Contract（前提条件：模数范围、是否互质、输入上限），再读函数体。凡是本书在正文里写\textbf{必须先检查}的地方，模板里都有对应的 \texttt{assert} 或返回值；提交时要按题面删掉 \texttt{assert} 以免被卡常数。
\end{tip}

完整 C++17 实现在资料包 \texttt{code/number\_theory.hpp}。本附录仅摘录关键部分，依赖同文件中的类型别名、\texttt{exgcd} 与相关函数；摘录不是可单独编译的提交文件。测试程序为 \texttt{code/selftest.cpp}。

\begin{longtable}{p{0.28\textwidth}p{0.61\textwidth}}
\toprule
模板 & 限制\\\midrule
\texttt{mul\_mod / pow\_mod} & 无符号 $64$ 位整数，模数必须大于零，乘积使用无符号 $128$ 位。\\
\texttt{inverse\_mod} & 不可逆时返回空 optional；模数一返回零是程序约定，不用于通常单位群讨论。\\
\texttt{exgcd} & 参数非负且不全为零；返回 gcd 及其一组整数线性组合系数。\\
\texttt{merge} & 输入及最终模数必须放入正的有符号 $64$ 位。合并最小公倍数超界时抛异常，不静默溢出。\\
\texttt{bsgs / exbsgs} & 返回最小非负解或空值；哈希表空间为平方根级，不能因为数值能放入 $64$ 位就对 $10^{18}$ 直接开表。\\
\texttt{PrimePowerBinom} & 素数参数需预先保证；默认最多建 $2\times10^6$ 个量级的表项。支持非负 $n,k$，$k>n$ 返回零。\\
\texttt{is\_prime / factor} & 判素适用于完整无符号 $64$ 位；分解使用可复现的默认随机种子，允许调用方修改，不是密码学实现。\\
\texttt{order} & 需要 $a$ 与 $m$ 互质，并提供正确的 $\varphi(m)$ 及其有序素因子列表。\\
\bottomrule
\end{longtable}

\section{安全模乘、快速幂与逆元}
\begin{lstlisting}
inline u64 mul_mod(u64 a, u64 b, u64 m) {
    assert(m > 0);
    return (u128)a * b % m;
}
inline u64 pow_mod(u64 a, u64 e, u64 m) {
    assert(m > 0);
    u64 r = 1 % m; a %= m;
    for (; e; e >>= 1, a = mul_mod(a, a, m))
        if (e & 1) r = mul_mod(r, a, m);
    return r;
}
inline std::optional<u64> inverse_mod(u64 a, u64 m) {
    assert(m > 0);
    // The residue for modulus 1 is returned as a programming convention.
    if (m == 1) return 0;
    i128 x, y;
    if (exgcd(a % m, m, x, y) != 1) return std::nullopt;
    return (u64)norm(x, m);
}
\end{lstlisting}

\section{最小非负离散对数}
\begin{lstlisting}
inline std::optional<u64> bsgs(u64 a, u64 b, u64 m) {
    assert(m > 0);
    if (m == 1) return 0;
    a %= m; b %= m;
    if (std::gcd(a,m) != 1) return std::nullopt;
    if (b == 1) return 0;
    if (std::gcd(b,m) != 1) return std::nullopt;
    u64 B = (u64)std::sqrt((long double)m);
    while ((u128)B*B < m) ++B;
    while (B && (u128)(B-1)*(B-1) >= m) --B;
    // Caller is responsible for the O(sqrt(m)) memory budget.
    std::unordered_map<u64,u64> baby;
    baby.reserve((std::size_t)(B*2+1));
    u64 v = 1;
    for (u64 j=0;j<B;++j) {
        baby.emplace(v,j); // Keep the earliest/smallest j.
        v = mul_mod(v,a,m);
    }
    u64 step = *inverse_mod(v,m), giant = b;
    for (u64 i=0;i<=B;++i) {
        auto it = baby.find(giant);
        if (it != baby.end()) {
            u128 x = (u128)i*B+it->second;
            if (x < m) return (u64)x;
        }
        giant = mul_mod(giant,step,m);
    }
    return std::nullopt;
}
inline std::optional<u64> exbsgs(u64 a, u64 b, u64 m) {
    assert(m > 0);
    if (m == 1) return 0;
    a %= m; b %= m;
    if (b == 1) return 0;
    if (a == 0) return b == 0 ? std::optional<u64>(1) : std::nullopt;
    u64 count = 0, k = 1;
    while (true) {
        u64 d = std::gcd(a,m);
        if (d == 1) break;
        if (b == k) return count;
        if (b % d) return std::nullopt;
        b /= d; m /= d; ++count;
        if (m == 1) return count;
        k = mul_mod(k,a/d,m);
    }
    auto inv = inverse_mod(k,m);
    assert(inv.has_value());
    auto tail = bsgs(a,mul_mod(b,*inv,m),m);
    if (!tail) return std::nullopt;
    return *tail + count;
}
\end{lstlisting}

\section{素数幂模数下的组合数}
\begin{lstlisting}
struct PrimePowerBinom {
    u64 p, q = 1;
    unsigned alpha;
    std::vector<u64> pref;
    // Contract: p is prime. Table size must be affordable.
    PrimePowerBinom(u64 prime, unsigned exponent,
                    u64 table_limit=2000000):p(prime),alpha(exponent) {
        assert(p>=2 && alpha>=1);
        for (unsigned i=0;i<alpha;++i) {
            if (q>table_limit/p) throw std::length_error("prime-power table too large");
            q*=p;
        }
        pref.assign(q+1,1);
        for (u64 i=1;i<=q;++i)
            pref[i]=(i%p) ? mul_mod(pref[i-1],i,q) : pref[i-1];
        assert(pref[q]==1 || pref[q]==q-1);
    }
    u64 valuation(u64 n) const {
        u64 e=0;
        while (n) { n/=p; e+=n; }
        return e;
    }
    u64 unit(u64 n) const {
        u64 r=1;
        while (n) {
            r=mul_mod(r,pref[n%q],q);
            if ((n/q)&1) r=mul_mod(r,pref[q],q);
            n/=p;
        }
        return r;
    }
    u64 choose(u64 n,u64 k) const {
        if (k>n) return 0;
        u64 e=valuation(n)-valuation(k)-valuation(n-k);
        if (e>=alpha) return 0;
        u64 r=unit(n);
        r=mul_mod(r,*inverse_mod(unit(k),q),q);
        r=mul_mod(r,*inverse_mod(unit(n-k),q),q);
        return mul_mod(r,pow_mod(p,e,q),q);
    }
};
\end{lstlisting}

\section{完整无符号 64 位的 Miller--Rabin}
\begin{lstlisting}
inline bool is_prime(u64 n) {
    if (n<2) return false;
    for (u64 p:{2ULL,3ULL,5ULL,7ULL,11ULL,13ULL,
                17ULL,19ULL,23ULL,29ULL,31ULL,37ULL})
        if (n%p==0) return n==p;
    u64 d=n-1; unsigned s=0;
    while (!(d&1)) { d>>=1; ++s; }
    for (u64 base:{2ULL,325ULL,9375ULL,28178ULL,
                   450775ULL,9780504ULL,1795265022ULL}) {
        u64 a=base%n;
        if (!a) continue;
        u64 x=pow_mod(a,d,n);
        if (x==1 || x==n-1) continue;
        bool pass=false;
        for (unsigned j=1;j<s;++j) {
            x=mul_mod(x,x,n);
            if (x==n-1) { pass=true; break; }
            if (x==1) break;
        }
        if (!pass) return false;
    }
    return true;
}
\end{lstlisting}

\section{怎样复现验证}

\begin{tip}[怎么做：对拍的最小可行版本]
这四条命令是\textbf{对拍}的最小可行版本：一份正确但慢的实现 + 一份你写的实现 + 随机数据。竞赛里 90\% 的\textbf{WA}能被它抓出来，包括取模负号、最小解、边界 $n=1$ 这类。
\end{tip}

在资料目录中运行：
\begin{lstlisting}[language=bash]
g++ -std=c++17 -O2 -Wall -Wextra -fsanitize=undefined \
    code/selftest.cpp -o nt_selftest
./nt_selftest
python validation/verify_math.py
python validation/check_happy.py
python code/vieta_construct.py --verify
\end{lstlisting}
C++ 测试程序以 Boost 的大整数作为组合数的独立对照；模板本身无需 Boost。Python 的数学验证仅依赖标准库。

\begin{tip}[提醒：Boost 只出现在自测里]
Boost 只在自测里当\textbf{独立裁判}（用 \texttt{cpp\_int} 精确算组合数再取模）。提交的代码只依赖 \texttt{number\_theory.hpp}。这个区分很重要：对拍用的实现要\textbf{笨且显然正确}，被测的实现要\textbf{快且接口干净}。
\end{tip}



\begin{tip}[代码与验证脚本的位置]
本附录摘录的模板，完整 C++17 实现在资料包的 \texttt{Number Theory/code/number\_theory.hpp}；
测试程序为 \texttt{Number Theory/code/selftest.cpp}；数学验证脚本为
\texttt{Number Theory/validation/verify\_math.py}、\texttt{check\_happy.py}
与 \texttt{Number Theory/code/vieta\_construct.py}。
这些文件相对本书源码目录的上一级，编译本书本身不需要它们。
\end{tip}

\chapter{教师使用：课堂设计、追问与评量}
\begin{chapterintro}[章节导读]
附录 A 解决“学生怎样练”，附录 B 解决“程序怎样写”；最后的附录 C 则把这两部分重新放回课堂。这里不新增一套数学主线，而是给出课程安排、追问、评分与验证框架，用来说明同一份数论材料怎样分别服务于讲授、训练和复习。
\end{chapterintro}

\section{如何安排一轮而不是一节“数论课”}
建议把本资料看作课程资源库，而不是一次讲完的幻灯片。每讲至少留下一个可验证的学习产出：一段具有输入输出契约的代码、一份没有跳步的推导、一组能击穿错误解法的反例。下面给出六个 $90$ 分钟课堂包；高级章节可另设专题阅读。

\begin{longtable}{p{0.15\textwidth}p{0.38\textwidth}p{0.37\textwidth}}
\toprule
阶段 & 通常的时间分配 & 教师的观察对象\\\midrule
诊断与情境 & $0$--$10$ 分钟：前置小测、具体难题 & 学生卡在数学概念，还是只忘了实现？\\
提出结构 & $10$--$30$ 分钟：换元、等价条件、不变量 & 能否说出“为什么这样设”？\\
证明与手算 & $30$--$50$ 分钟：核心证明、小数值验算 & 每一步是否依赖互质、素数、正性？\\
引导练习 & $50$--$70$ 分钟：同伴推导与纠错 & 先让学生指出错误，再补正确算法。\\
实现与迁移 & $70$--$85$ 分钟：代码接口、变式 & 能否从结论转换为边界安全的实现？\\
离堂检查 & $85$--$90$ 分钟：一道短题、一句解释 & 不只收答案，应收关键条件与理由。\\
\bottomrule
\end{longtable}

\section{课堂包一：exgcd 从“能解”到“求最小解”}
\paragraph{学习目标.} 学生能区分整数解与正整数解；能从一组特解生成全部解；能把两个正性约束转成同一个整数参数区间。
\paragraph{课前诊断.} 让学生回答 $6x+10y=7$ 为什么无解、$6x+10y=2$ 为什么可能有整数解却没有正整数解。不必先出现代码。
\paragraph{板书主线.}
\[
\text{公约数集合不变}\ \longrightarrow\ au+bv=d
\ \longrightarrow\ (x_0,y_0)=\frac cd(u,v)
\ \longrightarrow\ t\text{ 的区间}.
\]
\paragraph{引导顺序.}
先用 $252,198$ 手算 gcd；把余数倒代为整数线性组合。让学生自己发现乘 $c/d$，再用两组解相减证明 $\Delta x$ 必须是 $b/d$ 的倍数。最后讲 P5656，不要从通解公式直接开始背诵。
\paragraph{预设错误与回应.}
学生说“有解当且仅当 $c\mid d$”时，要求把 $a=6,b=10,c=32$ 代入判断，再找出 $(2,2)$；学生说“求到一个负数解说明无解”时，要求展示参数平移；学生把 ceil 写成 C++ 除法时，使用 $-5/3$ 进行对照。
\paragraph{同伴活动.} 一人负责给出一个特解，另一人只利用通解在区间内数解，再交换检查。两人使用不同特解时，最终正整数解集合应相同；这可以检验“特解选择不影响答案”的理解。
\paragraph{离堂题与作业.} 离堂做 A4；作业做 A1、A2、A8，再实现 P5656 的数学核心。评阅时要求附至少一个“有整数解但无正整数解”的测试。

\section{课堂包二：指数取模的三个不同问题}
\paragraph{学习目标.} 区分“算 $a^E$ 的余数”“求指数 $x$”“求最小周期”三种问题；知道欧拉函数是可用周期但不一定最小；知道非互质指数不能只保留余数。
\paragraph{课前诊断.} 写出模八的 $2^x$ 前五项，模七的 $2^x$ 前七项。请学生观察一个有尾巴、一个从开头就循环。
\paragraph{板书分栏.}
\begin{center}
\begin{tabular}{lll}
\toprule
问题 & 核心工具 & 不能丢的信息\\\midrule
计算幂余数 & 快速幂、扩展欧拉 & 指数是否达到阈值\\
求最小指数 & BSGS/exBSGS & 零指数、最小命中方式\\
求周期 & 阶与素因子试除 & 单位条件、正整数最小值\\\bottomrule
\end{tabular}
\end{center}
\paragraph{证明重点.} 欧拉定理只讲清一次“乘法置换剩余系”，不反复背费马与欧拉两套相同证明。扩展欧拉在素数幂上分类，把“不互质”变成“足够多素因子时归零”。
\paragraph{引导练习.} 先计算 $2^4\bmod8$，再对照指数只取模四的错误结果；接着让学生设计二元组接口 $(E\bmod q,E\ge q)$。随后手算 $2^x\equiv5\pmod{13}$，画出 baby/giant 的匹配表。
\paragraph{预设错误与回应.} “两个连续塔高余数相同就稳定”需要完整状态论证；“固定模数就能复用 BSGS 表”还要检查底数；“求阶时 BSGS 返回零”说明忘记阶要求最小\emph{正}指数。
\paragraph{离堂题与作业.} 做 B3、B4，要求先写能否使用阶的判断。提高作业为 P6736：只实现截断塔和模塔两个函数，再解释它们为什么返回不同信息，暂不要求最优时限提交。

\section{课堂包三：组合数中不能取逆的部分去哪了？}
\paragraph{学习目标.} 将组合数分为素数指数与单位部分；能解释 Lucas 只适合素数模数；能独立完成一个小合数模数的组合数计算。
\paragraph{情境.} 给出 $\binom{10}{4}\bmod12$，先允许学生尝试阶乘逆元法。等有人发现分母不能求逆时，再提出“并非不能除，而是不能在取模后直接除”。
\paragraph{板书主线.}
\[
 n!\ \longrightarrow\ p^{\nu_p(n!)}U_p(n)
 \quad\longrightarrow\ p^e\cdot\text{单位分式}
 \quad\longrightarrow\text{各素数幂余数}\quad\longrightarrow\mathrm{CRT}.
\]
\paragraph{证明顺序.} 先用倍数计数讲 Legendre；再用数位和讲 Kummer；最后讲 Lucas 的生成函数证明。扩展 Lucas 不是在原 Lucas 公式中把素数改成合数，而是更换为素数幂单位阶乘路线。
\paragraph{操作活动.} 将 $1,\ldots,10$ 中每个数写成二的幂乘奇数，排成两行。让学生亲自看见“删掉偶数”与“除去二的因子”得到不同乘积，然后推 $U_p(n)$ 的递归。
\paragraph{预设错误与回应.} $e\ge\alpha$ 可直接判零，但 $e<\alpha$ 时单位部分才决定剩余答案。单独的 $k!\bmod M=0$ 并不能说明组合数为零。若模数为一，整个计算都可提前结束。
\paragraph{离堂题与作业.} 离堂做 B9；基础作业做 B7、B8，提高作业做 P2183 和 P7488。要求附上一个最终余数为零但计数对象确实存在的例子。

\section{课堂包四：反演是为了分离变量}
\paragraph{学习目标.} 学生能将 $C(\gcd(i,j))$ 写成约数贡献和；能明确交换求和后每个因子的上界；能从反演过渡到前缀和递归。
\paragraph{情境.} 先画 $3\times3$ 的 gcd 矩阵并算总和十二。提出：“直接处理每个点，需要平方次；能否让每个公约数统一贡献？”
\paragraph{板书主线.}
\[
\gcd(i,j)=\sum_{d\mid i,d\mid j}\varphi(d)
\quad\Longrightarrow\quad
\sum_d\varphi(d)\floor{n/d}\floor{m/d}.
\]
再把 $\varphi$ 改为 $D=C*\mu$，得到一般加权公式。最后从 $\varphi*\one=\id$ 推导杜教筛，使同一卷积关系既用于拆点值，也用于拆前缀。
\paragraph{同伴活动.} 每组在一张纸上只写求和指标的条件，不写函数值：$1\le i\le n,1\le j\le m,d\mid i,d\mid j$。另一组把它改成先枚举 $d$ 的形式。这样可以暴露上界与整除条件的丢失，而不是让错误藏在长公式中。
\paragraph{预设错误与回应.} 把 gcd 拆成约数和后不能继续乘原 gcd；把 $i=da,j=db$ 后要同步修改两边上界；分块时两个商变化不同，必须取共同右端点。$S_\mu(n)$ 为零不代表它未被计算。
\paragraph{离堂题与作业.} 做 C2、C3；上机先写朴素 gcd 双循环和反演公式对拍，再实现记忆化杜教筛。提高学生可分析 P4240 的变长表项数。

\section{课堂包五：PN 与二维积性函数的“消主项”}
\paragraph{学习目标.} 理解为什么素数处匹配会让卷积商的支撑稀疏；理解二维坐标轴与对角线是两个不同层次的主要贡献；能正确叠加预处理和查询成本。
\paragraph{前置要求.} 必须已经熟练进行卷积的素数幂递推，知道形式幂级数常数项为一时可逆。若学生尚不能算出 C5 的三个局部系数，不宜直接讲 UOJ 885。
\paragraph{板书主线.}
\[
 f(p)=g(p)\Rightarrow h(p)=0\Rightarrow\text{PN 支撑};
\]
\[
 f(p^a,1)=g(p^a,1)\Rightarrow\text{去掉坐标轴};
\]
\[
 h=g_3*h'\Rightarrow h'(p^k,p^k)=0\Rightarrow\text{再去掉对角线}.
\]
\paragraph{引导练习.} 先做平方自由数计数，看到 $1-z^2$；再做 D4，看到 $1-XY$。这两个可完全手算的局部级数，是一般稀疏支撑论证的最好入口。
\paragraph{证明分层.} 基础提高班证明 PN 只有 $O(\sqrt n)$ 个；进一步证明同质因子对数为 $O(n)$；专题班再用收敛的局部乘积证明 $O((AB)^{1/3})$。不要同时把三层计数压缩成一句“显然很少”。
\paragraph{离堂题与作业.} 要求学生写出一个函数匹配素数处但 $h(p^3)\ne f(p^3)-g(p^3)$ 的例子，并指出总复杂度中哪项是前缀表预处理。研究作业为算法三中 $\rho<2/3$ 的收敛阈值，不直接要求完整最优实现。

\section{课堂包六：概率算法的答案为什么可信？}
\paragraph{学习目标.} 区分单边错误与零错误重试；知道概率的随机对象是底数而不是输入；能通过非平凡平方根解释判素和分解的联系。
\paragraph{情境.} 先用 $561$ 展示费马测试的失败，再展示 $263,166,67,1$ 的平方链。让学生自己指出 $67$ 为什么在素数模数下不可能出现。
\paragraph{板书对照.}
\begin{center}
\begin{tabular}{lll}
\toprule
工具 & 成功得到什么 & 失败时意味着什么\\\midrule
Miller--Rabin 见证 & 确定合数 & 未见证不自动证明为素数\\
Pollard--Rho 真因子 & 可直接验算的因子 & 更换参数重试\\
已知群阶倍数分裂 & 非平凡平方根 & 若全为 $\pm1$，重新随机\\\bottomrule
\end{tabular}
\end{center}
\paragraph{课堂实验.} 固定一个合数，枚举所有小底数，统计哪些会通过；再换合数，观察比例。必须先固定输入再讨论底数比例。让学生将输出因子乘回原数，理解可验证结果与概率判定结果的不同。
\paragraph{预设错误与回应.} “做了十次就绝对正确”需要明确随机误差；“这些七个底数适用任意大整数”要指出 $64$ 位范围；“得到 gcd 等于 $n$ 就递归 $n$”会无限循环；“发现 $-1$ 就得到真因子”用 gcd 立即反驳。
\paragraph{离堂题与作业.} 做 D1、D5。上机只对受控大小整数实现判素与分解，记录固定种子以便复现。进阶阅读 QOJ 1860 时，要求专门证明条件于 $a^T=1$ 的子群分布，而不是偷用 $\varphi(m)\mid T$。

\section{一次完整作业的评分框架}
\begin{longtable}{p{0.19\textwidth}p{0.12\textwidth}p{0.58\textwidth}}
\toprule
项目 & 建议分值 & 应检查的内容\\\midrule
建模与条件 & $20$ & 变量范围、互质／素数条件、是否允许零指数、对象是否有序。\\
正确性推导 & $30$ & 必要性与充分性、无漏解、无重复计数、不变量、递归终止。\\
算法与实现 & $20$ & 数据结构与推导一致，表的合法索引，模数一与负余数。\\
复杂度 & $15$ & 分解、预处理、每问及空间分开写；随机与最坏界不混淆。\\
验证与反例 & $15$ & 独立朴素程序、小范围穷举、边界、输出回代、失败案例。\\\bottomrule
\end{longtable}
建议允许学生用一个准确的反例修复先前不成立的结论，并对此给分。教师版的目标是建立可审查的推理，不是让学生把所有长公式背到最后。

\section{上机作业的最低测试清单}
\begin{enumerate}
\item \textbf{exgcd/CRT：} 不相容模数、一个模数整除另一个、余数为负、模数一、只有下界不限制余数、最小公倍数超出类型。
\item \textbf{离散对数：} 最小解为零、底数零、目标零、gcd 不为一、剥离后模数一、同余类重复导致的非最小命中。
\item \textbf{组合数：} $k=0,n$，$k>n$，$n=0$，$M=1$，模八的单位块乘积为一，赋值刚好等于模数指数。
\item \textbf{函数求和：} $n=1$，$S_\mu(n)=0$，两个商区间错位，多询问缓存，矩形两边差距极大。
\item \textbf{判素分解：} $0,1,2$，偶数、素数幂、Carmichael 数、强伪素数、接近 $2^{64}$ 的数、批量 gcd 等于 $n$ 的回退。
\item \textbf{构造：} 每组回代、组内大小范围、跨组全局去重、所有允许参数的规模验证。
\end{enumerate}

% \section{本次实际运行的检查}
% \begin{itemize}
% \item C++ 自测共执行 $1\,781\,581$ 次检查，覆盖带符号除法、逆元、最小离散对数、CRT、溢出拒绝、素数幂阶乘、组合数、判素、分解与求阶；启用了未定义行为检查器。
% \item 小范围 exgcd 正解区间与 P8255 构造实例共 $32\,165$ 项；LTE 有 $1088$ 项；P4240 覆盖 $40\times40$ 个矩形。
% \item 两阶段筛的 $720$ 个输入逐一核对了全部商集状态；二维卷积的两次分解用十二组系数与直接双循环对照。
% \item CF516E 与逐日模拟比较了 $14\,400$ 个小输入；AGC031F 的 $180$ 张小图对全部端点与余数询问比较完整状态图。
% \item WC2020 的 $62$ 个案例穷举子集覆盖，含两个原题样例；P7488 的递推与闭式在 $n\le12$ 上核对。
% \item P1400 覆盖 $k=2,\ldots,1000$，每个参数构造 $1000$ 组，验算方程、位数与全部 $3000$ 个数互异；逐参数 CSV 随附。
% \end{itemize}
% 这些记录说明具体实现和有限算例通过了检查，\textbf{不表示所有例题都已经在原评测平台提交满分，也不能替代无限范围的数学证明}。原题完整输入输出、极限常数、平台编译器与数据范围仍应在正式提交前核对。

% \chapter{主要勘误与严谨性补充}
% 下表记录合并时的实质性修正，既包括错误，也包括缺失条件与复杂度表述不完整。不是所有列出的原结论都错误：例如 Miller--Rabin 的 $1/2$ 界是可用弱界，但原证明需要重写；经验性构造也应与一般定理区分。

% \begingroup\small
% \begin{longtable}{p{0.055\textwidth}p{0.20\textwidth}p{0.645\textwidth}}
% \toprule
% 编号 & 位置 & 正确表述／本次处理\\\midrule
% \endfirsthead
% \toprule 编号 & 位置 & 正确表述／本次处理\\\midrule
% \endhead
% 1 & 第一稿：难度定位 & 基础整除可入门，但原根、二维积性函数及随机分解不能定位为小学二年级内容；改为按前置知识分层。\\
% 2 & 第一稿：P5656 & 整数解存在当且仅当 $\gcd(a,b)\mid c$，不是 $c\mid\gcd(a,b)$。\\
% 3 & 第一稿：P5656 & 从 $ax+by=d$ 的特解转到右边为 $c$ 时乘 $c/d$，不是 $d/c$。\\
% 4 & 第一稿：P5656 & 正整数解还需参数区间；有整数解不等于有正整数解；两个单独最小正值不一定能同时取得。\\
% 5 & 第一稿：P8255 & 应为 $d^2=\gcd(z/x,x^2)$，不是 $\gcd(z/x,a^2)$；还需 $x\mid z$ 与完全平方判定。\\
% 6 & 第一稿：欧拉展开 & $\gcd(i,n)$ 展开后约数项应为 $\varphi(d)$；保留原 gcd 会重复计数。\\
% 7 & 第一稿：CRT & 构造应为 $\sum_i a_i(M/m_i)((M/m_i)^{-1}\bmod m_i)$，不能把 $M/m_i$ 误写为 $m_i$。\\
% 8 & 两稿：CRT & 标准构造需模数两两互质；唯一性是模总乘积的唯一同余类。非互质时先检查相容条件，再取最小公倍数。\\
% 9 & 第一稿：屠龙勇士 & 选剑是“不高于生命值”的最大值，即 $\le$，不是严格小于。\\
% 10 & 第一稿：屠龙勇士 & 除同余式外还需攻击伤害下界 $x\ge\ceil{a_i/c_i}$；$a_i\le p_i$ 不是所有测试点的统一条件。\\
% 11 & 第一稿：LTE & $n=2^e u$ 去掉奇数部分后应为 $x^{2^e}-y^{2^e}$，不是 $x^{2^n}-y^{2^n}$。\\
% 12 & 第一稿：高阶箭号 & $3\uparrow^33$ 是高度 $3^{27}$ 的三的幂塔，不等于仅四层的 $3^{3^{27}}$。有限枚举需证明稳定高度。\\
% 13 & 第二稿：分配律 & $(a+b)c\equiv ac+bc$，不是 $ac+ab$。\\
% 14 & 第二稿：同余除法 & 同除公共因子后模数同步变为 $m/d$；等式中的商仍为原整数 $k$，不要求 $k/d$ 是整数。\\
% 15 & 第二稿：欧拉定理 & 模数必须为正，通常在 $m\ge2$ 讨论；“非负整数模数”会错误包含零。\\
% 16 & 第二稿：裴蜀证明 & 用“整数线性组合构成的子群”代替含糊的实向量空间说法，避免张成系数域不同导致误解。\\
% 17 & 第二稿：根数界 & 因式定理与非零根差的可逆性是在 $\mathbb F_p$ 中使用；合数模数可有超出次数的根数。\\
% 18 & 两稿：Wilson & 一般单位乘积中的符号按群结构分类；模二时 $1$ 与 $-1$ 是相同余数。\\
% 19 & 两稿：去因子阶乘 & $U_p(n)$ 是对每个数除尽素数因子后再相乘，不是直接删除所有该素数的倍数。\\
% 20 & 第二稿：扩展 Lucas & 查询成本含 $\log_p n$ 层递归与逆元成本，不能仅写成模数的函数；预处理要另计模数分解。\\
% 21 & 第一稿：P2183 & 不可行时原题输出 Impossible；可行方案数取模为零时仍应输出零，二者不等价。\\
% 22 & 第一稿：P7488 & 腰长彼此不同、底长彼此不同；不是全部腰长与底长数值整体两两不同，等边三角形允许出现。\\
% 23 & 第二稿：BSGS & 一般互质模数的指数上界来自欧拉定理，不来自只适用于素数的费马小定理。\\
% 24 & 第二稿：多询问 BSGS & 表可直接复用需要底数与模数都固定；仅固定模数不充分。\\
% 25 & 第二稿：BSGS 最小解 & 需规定重复余数存储策略与块枚举顺序；随便返回一次碰撞只保证一组解，不保证最小。\\
% 26 & 第二稿：exBSGS & 每层先检查零指数候选，除 gcd 前需 $d\mid b$；模数一、底数零要单独处理。\\
% 27 & 第二稿：原根判定 & 枚举的是 $\varphi(m)$ 的不同素因子，且须先保证 $g$ 为单位；不是任意整数因子。\\
% 28 & 第二稿：求阶复杂度 & 计算／分解 $\varphi(m)$ 是预处理成本，不能藏在每问快速幂成本中；求最小正阶不能直接接受指数零。\\
% 29 & 第二稿：CF516E & 同一循环只保留最小编号源的支配结论不成立；正文给反例并改成带权多源环扫线。\\
% 30 & 第二稿：CF571E & 参数是非负整数，不是“正整数且可为零”；各素因子的指数必须共享同一参数。\\
% 31 & 第二稿：WC2020 & 给出固定子集的最小覆盖询问数及其期望计数式，避免与在线未知答案的策略期望混淆。\\
% 32 & 第二稿：常见函数 & $\mu(p^k)=-[k=1]$；原稿已有误作 $I$ 的脚注，本次统一常一函数为 $\one$，约数个数为 $\tau$。\\
% 33 & 两稿：筛法与内存 & $O(N)$ 时间不等于小内存；$N=10^8$ 时多个 int 数组的空间可能不可接受，应按用途选择筛法。\\
% 34 & 第二稿：倍增估计 & 用端点代替一块各项是常数倍估计，应使用 $O,\Theta$ 或不等式，不是逐步精确相等。\\
% 35 & 第二稿：杜教筛 & 命中预处理范围应返回存储的前缀和，原稿的“返回 $\sqrt n$”是错误文字。补全阈值条件 $B\ge\sqrt n$。\\
% 36 & 第二稿：杜教筛 & $O(\sqrt n)$ 个状态，每个状态还有分块工作；未预处理总成本通常为 $O(n^{3/4})$，不是只数状态个数。\\
% 37 & 第二稿：PN 表示 & 不限制 $b$ 时 $a^2b^3$ 表示不唯一；要求 $b$ 平方自由可得唯一表示。计数上界不需要未经成立的双射。\\
% 38 & 第二稿：PN 总成本 & $O(\sqrt n)$ 是在相关前缀、局部系数与枚举准备完成后的支撑求和成本；完整算法还要加预处理。\\
% 39 & 第二稿：洲阁筛 & 第一阶段基函数必须能在非互质乘法中拆积，典型条件是完全积性；仅称积性不够。\\
% 40 & 第二稿：洲阁筛状态 & 明确排除整数一，原始前缀应减去一；素数幂枚举需有 $p^e\le v$ 上界；原地更新必须匹配旧层依赖。\\
% 41 & 第二稿：二维计数 & 从 $x,y\le n$ 放宽后应是 $xy\le n^2$，不是 $xy\le n$。\\
% 42 & 第二稿：二维支撑 & 对未出现的素数允许两边指数都为零；不能对所有素数都要求两边指数至少一。局部常数项必须为一。\\
% 43 & 第二稿：增长推广 & 应看第一个“非零”局部项而非“非负”项；本讲义明确归一化、非负整数系数与多项式增长条件。\\
% 44 & 第二稿：二维算法二 & 总成本必须加 $S,D$ 前缀预处理；算法三需要新的批量查询结构，不能把研究方向当成已完成实现。\\
% 45 & 第二稿：Carmichael & 费马性质针对 $\gcd(a,n)=1$ 的底数，而不是所有 $n\nmid a$ 的底数。\\
% 46 & 第二稿：Miller--Rabin & $1/2$ 可作为弱界；重新用局部循环群和 CRT 证明，并另给经典 $1/4$ 界。修正原证明中素数幂模数和指数的混写。\\
% 47 & 第二稿：一般分解 & 应说“没有已知的经典多项式时间算法”，不是断言不存在；复杂度变量是输入位数，并区别量子计算模型。\\
% 48 & 第二稿：Rho 两两比较 & $O(\sqrt p)$ 个状态若全部两两检查，会产生 $O(p)$ 次 gcd；判环才避免平方次检查。\\
% 49 & 第二稿：Rho 批处理 & $B\asymp\log n$ 时恢复项为 $O(\log^2n)$，不能不加条件略去；$n^{1/4}$ 是常用启发式量级而非无条件最坏界。\\
% 50 & 第二稿：$M=N\varphi(N)$ & $V=-1$ 时 $\gcd(V+1,m)=m$，必须找非平凡平方根。不能假设所有递归模数均满足 $\varphi(m)\mid M$，改用条件子群分析。\\
% 51 & 第二稿：韦达递降 & 另一根是整数但可能非正；应写清正整数解空间的停止条件。全局数字互异不同于三元组互异。\\
% 52 & 第二稿：P1400 范围 & 最多一百位指小于 $10^{100}$；经验性“128 位足够”不是普遍保证，平方验算还需更大范围。本次提供大整数实现与全参数有限验证。\\
% \bottomrule
% \end{longtable}
% \endgroup

% \section{排版与来源处理}
% 两稿统一为可用 XeLaTeX 编译的现代 UTF-8 中文文档，使用标准 $11$ 磅字号，避免不受支持的 $10.5$ 磅选项；统一定理、问题、解析、教学提示与警示样式。标题日期明确是本次整理日期，不冒充原讲稿日期；插图使用内嵌 TikZ，移除外部缺失图片依赖。


% \chapter{合并对应表、例题索引与延伸阅读}
% \section{知识点合并对应表}
% \begin{longtable}{p{0.34\textwidth}p{0.24\textwidth}p{0.31\textwidth}}
% \toprule
% 原稿内容 & 本书位置 & 整理方式\\\midrule
% 最大公约数、裴蜀、exgcd、步长循环 & 第 1 章 & 合并基本证明，补整数参数区间与边界。\\
% 同余、逆元、费马、欧拉、扩展欧拉 & 第 2 章 & 统一条件与记号，补截断指数接口与两种在线逆元方案。\\
% 中国剩余定理、扩展 CRT & 第 3 章 & 标准构造与逐式合并并列，补相容性、下界与溢出。\\
% 阶、原根、离散对数与高阶箭号 & 第 4 章 & 按周期到求指数重排，白泽教育归于此处而非 LTE。\\
% 第一稿的升幂引理 & 第 5 章 & 完整保留三类结论，重写证明，增加阶的提升例题。\\
% Wilson、阶乘、Legendre、Kummer、Lucas & 第 6 章 & 两稿合并，补单位阶乘与组合计数建模。\\
% 积性函数、卷积、线性筛、狄利克雷变换 & 第 7 章 & 统一 $\one,\mu,\tau$，区分完全积性与积性。\\
% 莫比乌斯／欧拉反演、gcd 求和 & 第 8 章 & 补通式、原题解析及多询问有效表存储。\\
% 倍增估计、块筛、杜教筛、PN & 第 9 章 & 补预处理／查询成本与局部支撑论证。\\
% 洲阁筛 & 第 10 章 & 写清两阶段状态、更新顺序与完全积性前提。\\
% 二维积性函数、UOJ 885 三类算法 & 第 11 章 & 补局部系数递推与更直接的支撑计数；最优改进列为研究拓展。\\
% 同余代数综合例题 & 第 12 章 & 补指数向量、环扫线、六状态图与期望计数解析。\\
% 素性判定、Miller--Rabin、Pollard--Rho & 第 13 章 & 区分概率界、确定性范围与启发式性能。\\
% 给定 $n,\varphi(n)$；构造 $n\varphi(n)$ & 第 14 章 & 补随机分裂的条件证明与选择性剪枝。\\
% 韦达定理、递降树、P1400 & 第 15 章 & 补正性边界、统一种子、全局去重与有限验证。\\
% 新增训练与教学资源 & 第 16 章及教师专章、附录 & 三十道分层练习、六个课堂包、模板与勘误。\\
% \bottomrule
% \end{longtable}

% \section{原有例题去重索引}
% \begin{longtable}{p{0.39\textwidth}p{0.12\textwidth}p{0.38\textwidth}}
% \toprule
% 题目 & 章节 & 教师版解析重点\\\midrule
% P5656 二元一次不定方程 & 1 & 通解、正解区间、无正解时的特殊输出。\\
% P8255 数学游戏 & 1 & 平方 gcd 的双向证明与解的唯一性。\\
% 五类乘法逆元问题 & 2 & 单次、线性、批量、减半递归与分块预处理。\\
% P4774 屠龙勇士 & 3 & 选剑、带系数同余、伤害下界、exCRT。\\
% P5605 小 A 与两位神仙 & 4 & 循环单位群、阶整除与非循环反例。\\
% P6736 白泽教育 & 4 & 幂塔截断、欧拉链稳定、三阶候选上界。\\
% 原稿两道赋值练习 & 5 & $7$-进差与二进偶指数 LTE。\\
% HDU 2973 YAPTCHA & 6 & 嵌套取整化为 Wilson 的判素示性函数。\\
% P2183 礼物 & 6 & 多项式系数与素数幂单位部分。\\
% P7488 黑色毛衣 & 6 & 阶梯棋盘递推与大参数合数取模。\\
% P2398 GCD SUM & 8 & 欧拉展开与交换求和。\\
% 一般加权 gcd 求和 & 8 & 任意数组的卷积差分与倍数后缀。\\
% P4240 毒瘤之神的考验 & 8 & $\mu^2/\varphi$、变长表与小／大参数分治。\\
% P4213 杜教筛 & 9 & 两个典型递归、记忆化与阈值分析。\\
% LOJ 6053 & 9 & 奇素数匹配 $\varphi$，单独处理素数二。\\
% P5325 & 10 & 素数贡献筛与按最小素因子补合数。\\
% UOJ 885 红场阅兵 & 11 & 坐标轴／对角线消去、支撑规模；算法三为研究拓展。\\
% CF516E & 12 & 带权多源环的前后缀最小值，修正原支配结论。\\
% CF571E & 12 & 两条指数向量射线的交集；大数据 Bonus 为研究拓展。\\
% AGC031F & 12 & 可逆状态、平移子群、每顶点六状态。\\
% P6730 WC2020 猜数游戏 & 12 & 阶分组、非单位短幂链、分量贡献期望。\\
% QOJ 4920 & 14 & 已知群阶倍数的随机非平凡平方根分裂。\\
% QOJ 1860 & 14 & 唯一性、选择性分解与条件子群分析。\\
% P1400 Easy Equation & 15 & 韦达跳跃、BFS、全局互异、大整数与全参数验证。\\
% \bottomrule
% \end{longtable}

% \section{来源与引用说明}
% \begin{enumerate}
% \item 原稿 A 为用户提供的《数论》文本；原稿 B 为用户提供、署名周雨衡、日期为 2025 年 9 月 26 日的同名讲义整理稿。本次不是对原作者原文的逐字复制，而是合并、重排、校订及新增教学阐释。
% \item 各原题的公开题面链接列在对应例题后，用于复查定义与输入输出。原稿提供的两个博客链接在本次访问时要求阅读密码，因此未把它们当成已阅读内容，也未尝试访问受保护正文；CF571E 与 AGC031F 的讲解根据数学推导及可公开核对的材料补充。
% \item AGC031F 的六状态路线可参考公开讲解：\url{https://www.luogu.com.cn/article/hz11gwzl}。本书补充了平移等价的条件、更新形式和商模数说明，并用完整小状态图对拍。
% \item UOJ 885 的延伸阅读为周康阳的专题文章：\url{https://www.cnblogs.com/zkyJuruo/p/18204106}。本书对算法一、二给出独立整理的卷积式与收敛支撑证明，算法三明确保留为专题研究方向。
% \item 组合数素数幂模数下更深入的算法可参考原稿链接：\url{https://faqak.github.io/misc/binom-mod-pe/}。本书采用适用于素数幂模数较小时的可解释预处理方案，不宣称已覆盖所有最优算法。
% \item 判素与分解可进一步阅读：\url{https://cp-algorithms.com/algebra/primality_tests.html} 与 \url{https://cp-algorithms.com/algebra/factorization.html}。固定底数保证、随机底数误差和 Rho 的启发式性能需分别理解。
% \end{enumerate}

% \section{后续修订建议}
% 对于新的题目或更大的数据范围，先重新写出函数接口、模数条件与资源上界，再复用本书的数学结构。尤其是多次询问、任意模数、高精度输入与最小解要求，经常改变一个本来正确模板的适用范围。任何新增推导都建议配一份独立朴素实现或最小反例，不把一次通过测试当成无条件证明。


% --------------------------------------------------------------------------
% 附录 D：来自《博弈论》教师版的“核心判定模板与接口”。
% --------------------------------------------------------------------------
\chapter{核心判定模板与接口}
\section{模板使用说明}
本附录摘录全书反复出现的判定框架，完整 C++17 实现见资料包 \texttt{code/selftest.cpp}；摘录是教学模板，不是可直接提交的完整程序。所有模板共用一个前提：游戏在有限步内结束、无法行动者负（三分类模板额外支持平局）。自测程序对有限范围的随机与穷举数据，把每个公式与暴力搜索逐点比对。

\begin{lstlisting}[language=C++]
// mex：集合中未出现的最小非负整数。
int mex(const std::vector<int>& s) {
    std::vector<char> seen(s.size() + 1, 0);
    for (int x : s)
        if (0 <= x && x < (int)seen.size()) seen[x] = 1;
    int m = 0;
    while (seen[m]) ++m;
    return m;
}
\end{lstlisting}

\begin{lstlisting}[language=C++]
// 模板一：DAG 上的记忆化 SG（隐式转移版）。
// 约定：move(u) 枚举 u 的全部后继；图无环（游戏有限步结束）。
// g(u) == 0 恰为 P 态（定理 5.x）。
struct SGTable {
    std::function<void(int, const std::function<void(int)>&)> move;
    std::vector<int> sg;
    std::vector<char> done;
    SGTable(int n, decltype(move) m) : move(m), sg(n, 0), done(n, 0) {}
    int dfs(int u) {
        if (done[u]) return sg[u];
        done[u] = 1;
        std::vector<int> nxt;
        move(u, [&](int v) { nxt.push_back(dfs(v)); });
        return sg[u] = mex(nxt);
    }
};
\end{lstlisting}

\begin{lstlisting}[language=C++]
// 模板二：一般状态图（允许有环）上的三分类。
// 返回 col：0 = P 态，1 = N 态，2 = 平局（过河卒框架）。
// u 是 N 态 <=> 存在 P 态后继；u 是 P 态 <=> 全部后继都是 N 态。
std::vector<int> solve_game(int n,
                            const std::vector<std::pair<int,int>>& edges) {
    std::vector<std::vector<int>> g(n), rg(n);
    for (auto [u, v] : edges) { g[u].push_back(v); rg[v].push_back(u); }
    std::vector<int> col(n, 2), und(n);
    std::vector<char> done(n, 0);
    std::queue<int> q;
    for (int u = 0; u < n; ++u) {
        und[u] = (int)g[u].size();
        if (und[u] == 0) { col[u] = 0; done[u] = 1; q.push(u); }
    }
    while (!q.empty()) {
        int u = q.front(); q.pop();
        for (int p : rg[u]) {
            if (done[p]) continue;
            if (col[u] == 0) { col[p] = 1; done[p] = 1; q.push(p); }
            else if (--und[p] == 0) { col[p] = 0; done[p] = 1; q.push(p); }
        }
    }
    return col;  // 仍为 2 的点：双方都不离开未定子图 => 平局。
}
\end{lstlisting}

\begin{lstlisting}[language=C++]
// 模板三：阶梯 Nim 判定（定理 2.x）。
// 输入 a[1..n]：第 i 级台阶上的石子数；a[0] 弃用（地面）。
// 每次把第 i 级的至少一颗石子移到第 i-1 级，不能行动者负。
// 先手必胜 <=> 奇数级异或和不为 0（只有奇数级参与！）
bool stair_nim_wins(const std::vector<long long>& a) {
    long long s = 0;
    for (int i = 1; i < (int)a.size(); i += 2) s ^= a[i];
    return s != 0;
}
\end{lstlisting}

Nim、反 Nim、Nim-k 与反 Nim-k 的判定公式各只有几行，不再单独摘录：它们在自测程序中以函数形式出现，并直接与暴力搜索对拍。

\begin{teach}[备课验证]
例题 [CF98E] 与 [THUPC 2023 Pre] 的 $2\times2$ 矩阵建议课前各解一遍数值，确认交点在 $0$ 到 $1$ 之间；本章练习 1--6 的数值答案（$p^{*}=\tfrac25,v=\tfrac15$；$(\tfrac14,\tfrac12,\tfrac14)$；$p^{*}=\tfrac23,q^{*}=\tfrac12$；闭式公式；$v=2$；$p^{*}=0,v=1$；扑克 $a=1$：$c=\tfrac13,p=\tfrac23,v=-\tfrac19$，$a=\tfrac12$：$\tfrac15,\tfrac45,-\tfrac1{25}$）应课前亲算一遍。附录自测程序应在课前编译运行一次，正常输出为一万二千余项检查全部通过。若你修改了任何判定模板，先跑自测再进课堂。
\end{teach}
\section{怎样复现验证}
在资料目录中运行：
\begin{lstlisting}[language=bash]
g++ -std=c++17 -O2 -Wall -Wextra code/selftest.cpp -o gt_selftest
./gt_selftest
\end{lstlisting}
自测程序完成五组有限范围验证：(1) “每次取 $1$ 或 $2$ 颗”的 SG 值周期 $3$；(2) Nim、反 Nim、Nim-k（$k=1,2$）与反 Nim-k（$k=1,2$）的判定公式对全部 $4$ 堆以内、每堆不超过 $5$ 颗的局面与暴力博弈搜索一致；(3) 三分类模板与不动点迭代在随机带环图上一致；(4) 阶梯 Nim 判定与暴力搜索一致；(5) 记忆化 SG 与显式转移表一致。暴力验证依赖的是“状态数小到可以穷举”这一事实，与正文证明互为印证而非替代；更大的参数范围仍以正文证明为准。


\begin{tip}[代码与验证脚本的位置]
本附录摘录的判定模板，完整 C++17 实现在资料包的 \texttt{Game Theory/code/selftest.cpp}。
自测程序对有限范围的随机与穷举数据，把每个公式与暴力搜索逐点比对；
它相对本书源码目录的上一级，编译本书本身不需要它。
\end{tip}
